Este informe compara la teoría con la práctica en dos problemas
clásicos: ordenar un arreglo de enteros y multiplicar matrices
cuadradas. Del lado del ordenamiento se implementaron Merge Sort,
Quick Sort, Patience Sort y Sort; del lado de matrices,
el método Naive y el de Strassen. Todos tienen una cota de complejidad
bien conocida, pero eso no garantiza que se comporten como uno espera
al correrlos de verdad: cosas como el tamaño de la entrada, qué tan
ordenados o repetidos están los datos, o simplemente cómo está escrito
el código, terminan pesando tanto como la notación asintótica misma.
Ese es justamente el punto de este trabajo: medir tiempo y memoria en
distintos escenarios y ver dónde la teoría predice bien y dónde no.